
\documentclass{article}
\usepackage{graphicx} % Required for inserting images

\usepackage{amsmath}
\usepackage{amsfonts}
\usepackage{tikz-cd}

\title{test}
\date{May 2025}

\begin{document}

\textbf{Exercise 2.3}

\vspace{0.5cm}

if $\lambda \not\in \sigma(a^*)$ then there is some $b\in \mathcal A$ with $$b(a^*- \lambda) = (a^*- \lambda)b = 1$$, then $$(a- \bar\lambda)b^* = b^*(a- \bar\lambda) = 1 $$ and $\lambda \not \in \overline{\sigma(a)}$ so then $$\overline{\sigma(a)} \subseteq \sigma(a^*)$$ and this is valid if we replace $a$ by $a^*$ aswell: $$\overline{\sigma(a^*)} \subseteq \sigma(a)$$ and then take conjugate on both sides to obtain reverse inclusion $$\sigma(a^*) \subseteq \overline{\sigma(a)}$$ which establishes $\overline{\sigma(a)} = \sigma(a^*)$.

\vspace{0.5cm}

\textbf{Exercise 2.4} 

\vspace{0.5cm}

$ea = a$ implies $a^*e^* = a^*$ and $a \mapsto a^*$ is surjective so that $ae^* = a$ for all $a \in \mathcal A$ meaning $e^*$ is a right unit, then $ee^* = e$ since $e^*$ is a right unit and $ee^* = e^*$ since $e$ is a left unit, then $e = e^*$ and $e$ is a unit.

\vspace{0.5cm}

\textbf{Exercise 2.5}

\vspace{0.5cm}

$x = L, \ y = R \in \mathcal B(l^2(\mathbb N))$, left and right shift operator in the Banach algebra of bounded operators on $l^2(\mathbb N)$

\vspace{0.5cm}

\textbf{Exercise 2.6}

\vspace{0.5cm}

$f = \delta_1$ works. $\delta_{1k} = 1$ if $k = 1$ and $= 0$ if $k \neq 1$. $\mathcal A =l^1(\mathbb Z)$ is a commutative unital Banach algebra so then $\sigma(f) = \hat f(\Delta_\mathcal A)$ (Theorem 2.5.2) and we see in exercise 2.16 in detail why $\hat f(\Delta_\mathcal A) = S^1$ so $\sigma(f) = S^1 \not \subseteq \mathbb R$.

\vspace{0.5cm}

\textbf{Exercise 2.9}

\vspace{0.5cm}

$r(a) = \lim\limits_{n \rightarrow \infty} ||a^n||^{1/n}$ and $\mathcal B$ becomes a banach algebra with norm from $\mathcal A$ so we can assume their norms agree $|| x||_\mathcal A = ||x||_\mathcal B$. $r_\mathcal A(x) = \lim\limits_{n \rightarrow \infty} ||a^n||_\mathcal A^{1/n} = \lim\limits_{n \rightarrow \infty} ||a^n||_\mathcal B^{1/n} = r_\mathcal B(x)$

\vspace{0.5cm}

\textbf{Exercise 2.10}

\vspace{0.5cm}

$\mathcal B = B(x,y,1)$ (Banach algebra generated by x, y and 1) its unital and commutative, so we can apply the identity $||\hat b||_\infty = r_\mathcal{B}(b)$ for all $b \in \mathcal B$ where $\hat b : \Delta_{\mathcal B} \rightarrow \mathbb{C}$ (Theorem 2.5.2 (c)). Then we see that $r_\mathcal B(x+y) \leq r_\mathcal B(x) + r_\mathcal B(y)$ from triangle inequality of sup norm and also $r_\mathcal B(xy) \leq r_\mathcal B(x)r_\mathcal B(y)$ from basic sup norm thing again.

\vspace{0.5cm}

\textbf{Exercise 2.11}

\vspace{0.5cm}

Pass to the unitization first, since the inclusion $\mathcal A \rightarrow \mathcal  A^e$ is a bounded map then if we prove the statement for unital Banach algebras then we prove it for non-unital aswell since $1 \leq ||e-f||_{\mathcal A^e} \leq ||e-f||_\mathcal A$. 

So assume $\mathcal A$ is unital and let $\mathcal B =B(e,f,1)$ be the unital Banach algebra generated by $e \ \& \ f$ inside $\mathcal A$. The spectral mapping theorem for polynomials says that in a unital Banach algebra we have $p(\sigma(x)) = \sigma(p(x))$ for a polynomial $p$ with complex coefficients. Set $x = e-f$ then let $p(z) = z^3 - z$ and we see that $p(x) =x^3 - x = 0$ can be checked. So then $\sigma(p(x)) = \{0\}$ which means that all elements of $\sigma(x)$ have to be roots of $p$ by spectral mapping theorem. $\sigma(x) \subseteq \{0,-1,1\}$, the only thing we have to rule out is that $\sigma(x) = \{0\}$ because if this isnt true then we get spectral radius $r(x) = 1$ and so $||e-f|| = ||x|| \geq r(x) = 1$ and we are done. So assume $\sigma(x) = \{0\}$ then $\sigma(x-1) = \{-1\}$ and $\sigma(x+1) = \{1\}$ which both dont include $0$ which means that $x-1$ and $x+1$ are both invertible, this would be a contradiction since $x^3 - x = x(x-1)(x+1) = 0$ and if we multiply by the inverse of $x-1$ and $x+1$ we get $x = 0$ i.e $e = f$ which we assume was not true, contradiction.

\vspace{0.5cm}

\textbf{Exercise 2.12}

\vspace{0.5cm}

Let $f \in \mathcal A$, then $f$ is uniformly continuous so if $\epsilon > 0$ and $\delta > 0$ exists by uniform continuity of $f$ and we pick $r < 1$ with $|r-1| < \delta$ we can guarantee that $||f_r - f|| < \epsilon$ where $f_r(z) := f(rz)$ since $|zr - z| \leq |1-r| < \delta$. $f_r$ is holomorphic on the interior of the disk with radius $1/r > 1$. Then by complex analysis its power series converges uniformly to it on the closed disk of radius $1$. This means that there is a polynomial $p(z)$ with $||f_r-p|| < \epsilon$ and thus 

$$||f - p|| \leq ||f - f_r|| + ||f_r - p|| < 2 \epsilon$$

\vspace{0.5cm}

\textbf{Exercise 2.13}

\vspace{0.5cm}

$f(z) = z + i$ works as a counterexample to the $C^*-$identity $||ff^*|| = ||f||^2$ since $||f|| = 2$ and $f(z)f^*(z) = z^2 + 1$ so $||ff^*|| = 2$, the $C^*$ identity would then be $2 = 4$.

\vspace{0.5cm}

\textbf{Exercise 2.14}

\vspace{0.5cm}

First assume $X$ is compact. Then following the hint of the exercise (by contradiction) if $I$ is a proper ideal with the assumption $\forall x \exists f \in I : f(x) \neq 0$ then $U_f = f^{-1}(\mathbb{R} - \{0\})$ cover $X$ and by compactness there are finitely many $f_k \in I$ so that $U_{f_k}$ cover $X$. Since $I$ is an ideal we have $\bar f_k f_k = |f_k|^2 \in I$ aswell and then also $f(x) = \sum\limits_{k} |f_k(x)|^2$ is in $I$ and is never zero. So therefore its an invertible element of the algebra and it cant be contained in a proper ideal (since then $1 = \frac{1}{f} f$ would also be in $I$ and then $\forall g:g = g 1 \in I$ so $I = C(X)$ and its not a proper ideal). 

\vspace{0.2cm}

So then if $m$ is a maximal ideal in $C(X)$ there is an $x \in X$ with $m \subseteq \ker(m_x)$ and by maximality of $m$ then $m = m_x$. We see then that $X \rightarrow \Delta_\mathcal A$ is surjective. Its also injective because points can be separated by real valued functions in a locally compact Hausdorff space. $X \rightarrow \Delta_\mathcal A$ is continuous because if $x_i \rightarrow x$ then for each continuous $f: X \rightarrow \mathbb C$, $m_{x_i}(f) \rightarrow m_x(f)$ is equivalent to $f(x_i) \rightarrow f(x)$ so $m_{x_i} \rightarrow m_x$ in the weak-$*$ topology. Since a continuous bijection $K \rightarrow H$ where $K$ is compact and $H$ is Hausdorff is a homeomorphism we are done.

\vspace{0.2cm}

Now lets deal with noncompact $X$.

\vspace{0.1cm}

\begin{center}
\begin{tikzcd}
X \arrow[r] \arrow[d] & \Delta_{C_0(X)} \arrow[d] \arrow[r, "\text{compactification}"] & \Delta_{C_0(X)^e} \arrow[ld, "\approx"] \\
X^\infty \arrow[r, "\approx"]    & \Delta_{C(X^\infty)}                                           &                                 
\end{tikzcd}
\end{center}

\vspace{0.2cm}

This diagram commutes, the arrow $\Delta_{C_0(X)^e} \rightarrow \Delta_{C(X^\infty)}$ is a homeomorphism and comes from lemma 2.3.2. If we remove the point at infinity of $X^\infty$ and $\Delta_{C(X^{\infty})}$ we get 

\begin{center}
\begin{tikzcd}
X \arrow[r] \arrow[d, "\approx"]           & \Delta_{C_0(X)} \arrow[d, "\approx"'] \\
X^\infty - \{\infty\} \arrow[r, "\approx"] & \Delta_{C(X^\infty)}- \{m_\infty\}   
\end{tikzcd}
\end{center}

where the vertical arrows become homeomorphisms since we saw they were both previously one-point compactifications.

\vspace{0.5cm}

\textbf{Exercise 2.16}

\vspace{0.5cm}

The idea is to use Theorem 2.5.2 $\sigma(a) = \hat a(\Delta_\mathcal A)$ to prove that $0 \not \in \sigma(a)$ which would mean that $a$ is invertible.

Let's first find all bounded multiplicative nonzero functionals $m:l^1(\mathbb Z) \rightarrow \mathbb C$. We start with the observation that the dual space of $l^1$ is $l^\infty$ so then $m$ is represented by $(...,m_{-1}, m_0, m_1,...) \in l^\infty(\mathbb Z)$ as $m(x) = \sum\limits_k x_k m_k$ then we look at the relation $m(x*y) = m(x)m(y)$ with $x = \delta_i$ and $y = \delta_j$ and obtain $m(\delta_i * \delta_j) = m(\delta_{i+j}) = m_{i+j}$ which is also equal to $m(\delta_i)m(\delta_j) = m_im_j$ so then $m_{i+j} = m_i m_j$ and so all of $m$ is determined by $m_1 \in \mathbb C$, we cant have $|m_1| \neq 1$ since then either $m_k = (m_1)^k$ or $m_{-k} = (m_1)^{-k}$ will grow unbounded in size as $k \geq 0$ grows larger if $|m_1| > 1$ or $|m_1| < 1$ respectively and then $m$ wouldn't be bounded so then $m_1 \in S^1$ ($m_1$ can also never be zero since then $m = 0$ which is a multiplicative functional but not nonzero). Now if $z \in S^1$ then $m(x) = \sum\limits_{k \in \mathbb Z} x_k z^k$ is a nonzero multiplicative functional so this assignment $z \rightarrow m$ gives a map $S^1 \rightarrow \Delta_\mathcal A$ that is surjective by what we just did, its also injective since we can recover $z$ from $m$ by $z = m(\delta_1)$. Its continuous because if $z_n \rightarrow z$ then $m_n(x) =\sum\limits_{k \in \mathbb Z} x_k z_n^k \rightarrow \sum\limits_{k \in \mathbb Z} x_k z^k = m(x)$ for each $x \in l^1(\mathbb Z)$, that's just continuity of the function $f(z) = \sum\limits_{k \in \mathbb Z} x_k z^k$ on $S^1$ and that's exactly weak-* convergence in $\Delta_\mathcal A$, $m_n \rightarrow m$ so then $S^1 \xrightarrow \approx \Delta_\mathcal A$ is continuous and this also a homeomorphism since its a bijection Compact $\rightarrow$ Hausdorff. We know that $\sigma(x) = \hat x(\Delta_\mathcal A)$ and under this homeomorphism we see that this equals the image of $S^1$ under $\sum\limits_{k \in \mathbb Z} x_k z^k$ so then finally if $f(x)$ and $a_n$ is as in the problem statement then $\sigma(a) = f(\mathbb R) \not \ni 0$ and $a$ is invertible.

\vspace{0.5cm}

\textbf{Exercise 2.17}

\vspace{0.5cm}

Let $\mathcal C = C^*(\text{im}(\phi)) \subseteq \mathcal B$ be the $C^*$ closure of the image of $\phi$ in $\mathcal B$. Then if $m : \mathcal C \rightarrow \mathbb C$ be a nonzero multiplicative linear functional then $m$ is nonzero on $\text{im}(\phi)$ so $m \circ \phi$ is nonzero and we get a continuous map $\phi^* : \Delta_\mathcal C \rightarrow \Delta_\mathcal A$ and the diagram

\begin{center}
\begin{tikzcd}
\mathcal A \arrow[r, "\phi"] \arrow[d, "\approx"] & \mathcal C \arrow[d, "\approx"] \\
C(\Delta_\mathcal A) \arrow[r, "{\phi^*}^*"]      & C(\Delta_\mathcal C)           
\end{tikzcd}
\end{center}

which proves $\phi$ is a bounded $*-$homomorphism since ${\phi^*}^*$ is.

\vspace{0.5cm}

\textbf{Exercise 2.18}

\vspace{0.5cm}

The functional calculus in the unital case gives us $\tilde\Phi : C(\sigma_{\mathcal A}(a)) \rightarrow \mathcal A^e$, next we look at the set $Q$ of functions of the form $q(x) = p(x,\overline x)$ where $p$ is a polynomial in two variables and has constant term $ = 0$. These functions $q \in Q$ are dense in $C_0(\sigma_{\mathcal A}(a))$, if $q$ is such a function then $\tilde\Phi(q) \in \mathcal A$ so $\tilde\Phi$ restricts to a map $Q \rightarrow \mathcal A$ and by taking closure $\tilde \Phi$ restricts to an isometry $\Phi:C_0(\sigma_{\mathcal A}(a)) \rightarrow \mathcal A$ since the inclusion $\mathcal A \subseteq \mathcal A^e$ is an isometry. Since the $C^*-$closure of $q(x) = x$ inside of $C_0(\sigma_{\mathcal A}(a))$ is the entirety of $C_0(\sigma_{\mathcal A}(a))$ then the $C^*-$closure of $a \in \mathcal A$ (which is $C^*(a)$) is the image of $C_0(\sigma_{\mathcal A}(a))$ under $\Phi$. I.e $\Phi(C_0(\sigma_\mathcal A(a))) = C^*(a)$. This covers existence.

\vspace{0.3cm}

Uniqueness: If we have two such $\Psi,\Phi : C_0(\sigma_\mathcal A(a)) \rightarrow \mathcal A$ then 
$$\Phi^e, \Psi^e : C_0(\sigma_\mathcal A(a))^e \rightarrow \mathcal A^e$$ 
are two unital $*-$homomorphisms and precomposing with the unital $*-$homomorphism 
$$C(\sigma_\mathcal A(a)) \rightarrow C_0(\sigma_\mathcal A(a))^e: f \mapsto (f - f(0), f(0))$$ 
we get two unital $*-$homomorphisms 
$$\Phi', \Psi' : C(\sigma_\mathcal A(a)) \rightarrow \mathcal A^e$$ 
with $\Phi'(\text{Id}) = \Psi'(\text{Id}) = a$ and so by theorem 2.7.3 they are equal which implies $\Phi = \Psi$.

\vspace{0.5cm}

\textbf{Exercise 2.19}

\medskip

\noindent\emph{Existence.} 

\smallskip

\noindent\textbf{Unital case.} Let \(\mathcal A\) be a unital \(C^*\)-algebra and let \(a\in\mathcal A\) be positive. Then \(\sigma(a)\subset[0,\infty)\), and the continuous functional calculus applied to the continuous function \(f(t)=t^{1/n}\) on \(\sigma(a)\) yields a positive element \(b=f(a)\in C^*(a,1)\) with
\[
b^n=a.
\]
Thus a positive \(n\)-th root exists and lies in the unital commutative subalgebra \(\mathcal C:=C^*(a,1)\).

\smallskip

\noindent\textbf{Non-unital case.} If \(\mathcal A\) is not unital, pass to the unitization \(\mathcal A^e\). The element \((a,0)\in\mathcal A^e\) is positive and has the same spectrum as \(a\) (by the definition of spectrum in the non-unital setting), so the unital functional calculus in \(\mathcal A^e\) produces a positive \(n\)-th root \((b,\lambda)\in\mathcal A^e\) with \((b,\lambda)^n=(a,0)\). Writing the product in coordinates gives \((b,\lambda)^n = (\text{stuff},\lambda^n)=(a,0)\), so \(\lambda^n=0\) and hence \(\lambda=0\). Therefore \(b^n=a\) in \(\mathcal A\). This completes the existence part.

\medskip

\noindent\emph{Uniqueness.}

\smallskip

\noindent\textbf{Unital case.} Again let $\mathcal A$ be unital and let \(b\in\mathcal C=C^*(a,1)\) be the positive \(n\)-th root obtained from functional calculus. Suppose \(d\in\mathcal A\) is another positive element with \(d^n=a\). Then \(d\) commutes with \(a\): indeed
\[
ad = d^n d = d d^n = da,
\]
so the subalgebra \(\mathcal D:=C^*(a,d,1)\) is a unital commutative \(C^*\)-subalgebra of \(\mathcal A\). Therefore since both $d$ and $b$ are positive $n-$th roots of $a$ in the $C^*$ algebra $\mathcal D = C(\Delta_\mathcal D)$ we are done since positive n-th roots are unique in the case of positive functions.

\smallskip

\noindent\textbf{Non-unital case.} If \(\mathcal A\) is non-unital and \(d,b\in\mathcal A\) are positive \(n\)-th roots of \(a\), then \((d,0)\) and \((b,0)\) are positive \(n\)-th roots of \((a,0)\) in the unitization \(\mathcal A^e\). By the uniqueness already proved in the unital case we have \((d,0)=(b,0)\), hence \(d=b\). This completes the uniqueness proof.

\medskip

\vspace{0.5cm}

\textbf{Exercise 2.20}

\vspace{0.5cm}

The idea of the proof is very similar to 2.19 where we reduced to the case of a $C^*-$algebra $C(X)$ for a compact topological space $X$.

First let's cover existence. Assume the result is true for the unital $C^*-$ algebra $\mathcal A^e$. So that $(a,0) =(a^+,\lambda) - (a^-,\lambda)$ for two positive elements of $A^e$ where $\lambda \in \mathbb C$ but from the condition $(a^+,\lambda) \cdot (a^-,\lambda) =(0,0) $ we see that $\lambda^2 = 0$ and so $\lambda = 0$ and $a = a^+ - a^-$ where $a^+,a^-$ are both positive and $a^+a^- = a^-a^+ = 0$ so we have reduced existence in non-unital case to existence in the unital case which we will cover next.

If $\mathcal A$ is a unital $C^*$-algebra the existence result follows from the functional calculus letting $f^+(x) = \max(x,0)$ and $f^-(x) = - \min(x,0)$ then $f^+(x) - f^-(x) = x$ and $f^-(x)f^+(x) = f^+(x)f^-(x) = 0$ and they are both positive so $a^+ = f^+(a)$ and $a^- = f^-(a)$ are two positive elements that satisfy our requirements. Also $a^+,a^- \in C^*(a,1)$ by functional calculus.

Now we prove uniqueness in the unital setting, if $a = d^+ - d^-$ where $d^+,d^-$ satisfy our requirements we proceed exactly analogously to exercise 2.19, let 
$$\mathcal D = C^*(a,d^+, d^-, 1)$$
It is commutative since $ad^+ = (d^+ - d^-)d^+ = (d^+)^2 = d^+ (d^+ - d^-) = d^+ a$ and similarly for $d^-$). Therefore since $\mathcal D = C(\Delta_\mathcal D)$ $a^+ = d^+$ and $a^- = d^-$ since positive-negative decomposition is unique in a $C^*-$algebra of continuous functions on a topological space.

Uniqueness for non-unital $C^*-$algebras follows aswell by considering the unitization $\mathcal A^e$.

\vspace{0.5cm}

\end{document}
